% GENERATED FILE. Edit canonical lessons, then run npm run generate:books.
% canonical-source-hash: fnv1a64:81a51e329e88acfb
% canonical-lessons: ML-C311-talotuka, ML-C311-otukkuka, ML-C311-polikkuka, ML-C311-kaykkuka, ML-C311-pataruka

\chapter{To stroke, To tidy away, To pull down, To bear fruit, To spread}
\label{ch:ml-to-stroke-to-tidy-away-to-pull-down-to-bear-fruit-to-spread}
\hlchaptermodality{311}

\begin{chapteropening}
\textbf{By the end of this chapter:} \emph{I can say to stroke, to tidy away, to pull down, to bear fruit and to spread.}

Complete the last lesson of chapter 311: I can say to stroke, to tidy away, to pull down, to bear fruit and to spread.
\end{chapteropening}

\section[\texorpdfstring{\ml{തലോടുക} (talōṭuka)}{തലോടുക (talōṭuka)}]{\ml{തലോടുക} (talōṭuka) — to stroke, to pat gently}

\label{lesson:ML-C311-talotuka}



\begin{quote}
Before the new one: say the Malayalam for to give up, then the Malayalam for to repeat.
\end{quote}

\subsection*{You'll want to know: \ml{തലോടുക}}
\textbf{\ml{തലോടുക}} — \emph{talōṭuka} — \textquotedblleft{}to stroke, to pat gently\textquotedblright{}.

\begin{grammarlens}[title={What you've built}]
One more word for this chapter.
\end{grammarlens}

\subsection*{Your turn}
\begin{itemize}
\raggedright
  \item \emph{Say it:} \emph{talōṭuka}
  \item \emph{Say it:} \emph{talōṭuka}, once more
  \item \emph{Say it:} say \emph{talōṭuka}
  \item \emph{Recall:} say the Malayalam for to give up, then the Malayalam for to repeat, then say \emph{talōṭuka} again
\end{itemize}

\begin{tcolorbox}[breakable,colback=teal!4,colframe=teal!35!black,title={Before you move on}]
What does \ml{തലോടുക} mean? (\textquotedblleft{}To stroke, to pat gently\textquotedblright{}.) Say it once more.
\end{tcolorbox}
\section[\texorpdfstring{\ml{ഒതുക്കുക} (otukkuka)}{ഒതുക്കുക (otukkuka)}]{\ml{ഒതുക്കുക} (otukkuka) — to tidy away, to put in order}

\label{lesson:ML-C311-otukkuka}



\begin{quote}
Before the new one: say the Malayalam for to repeat, then the Malayalam for to stroke.
\end{quote}

\subsection*{You'll want to know: \ml{ഒതുക്കുക}}
\textbf{\ml{ഒതുക്കുക}} — \emph{otukkuka} — \textquotedblleft{}to tidy away, to put in order\textquotedblright{}.

\begin{grammarlens}[title={What you've built}]
One more word for this chapter.
\end{grammarlens}

\subsection*{Your turn}
\begin{itemize}
\raggedright
  \item \emph{Say it:} \emph{otukkuka}
  \item \emph{Say it:} \emph{otukkuka}, once more
  \item \emph{Say it:} say \emph{otukkuka}
  \item \emph{Recall:} say the Malayalam for to repeat, then the Malayalam for to stroke, then say \emph{otukkuka} again
\end{itemize}

\begin{tcolorbox}[breakable,colback=teal!4,colframe=teal!35!black,title={Before you move on}]
What does \ml{ഒതുക്കുക} mean? (\textquotedblleft{}To tidy away, to put in order\textquotedblright{}.) Say it once more.
\end{tcolorbox}
\section[\texorpdfstring{\ml{പൊളിക്കുക} (poḷikkuka)}{പൊളിക്കുക (poḷikkuka)}]{\ml{പൊളിക്കുക} (poḷikkuka) — to pull down, to demolish}

\label{lesson:ML-C311-polikkuka}



\begin{quote}
Before the new one: say the Malayalam for to stroke, then the Malayalam for to tidy away.
\end{quote}

\subsection*{You'll want to know: \ml{പൊളിക്കുക}}
\textbf{\ml{പൊളിക്കുക}} — \emph{poḷikkuka} — \textquotedblleft{}to pull down, to demolish\textquotedblright{}.

\begin{grammarlens}[title={What you've built}]
One more word for this chapter.
\end{grammarlens}

\subsection*{Your turn}
\begin{itemize}
\raggedright
  \item \emph{Say it:} \emph{poḷikkuka}
  \item \emph{Say it:} \emph{poḷikkuka}, once more
  \item \emph{Say it:} say \emph{poḷikkuka}
  \item \emph{Recall:} say the Malayalam for to stroke, then the Malayalam for to tidy away, then say \emph{poḷikkuka} again
\end{itemize}

\begin{tcolorbox}[breakable,colback=teal!4,colframe=teal!35!black,title={Before you move on}]
What does \ml{പൊളിക്കുക} mean? (\textquotedblleft{}To pull down, to demolish\textquotedblright{}.) Say it once more.
\end{tcolorbox}
\section[\texorpdfstring{\ml{കായ്ക്കുക} (kāykkuka)}{കായ്ക്കുക (kāykkuka)}]{\ml{കായ്ക്കുക} (kāykkuka) — to bear fruit}

\label{lesson:ML-C311-kaykkuka}



\begin{quote}
Before the new one: say the Malayalam for to tidy away, then the Malayalam for to pull down.
\end{quote}

\subsection*{You'll want to know: \ml{കായ്ക്കുക}}
\textbf{\ml{കായ്ക്കുക}} — \emph{kāykkuka} — \textquotedblleft{}to bear fruit\textquotedblright{}.

\begin{grammarlens}[title={What you've built}]
One more word for this chapter.
\end{grammarlens}

\subsection*{Your turn}
\begin{itemize}
\raggedright
  \item \emph{Say it:} \emph{kāykkuka}
  \item \emph{Say it:} \emph{kāykkuka}, once more
  \item \emph{Say it:} say \emph{kāykkuka}
  \item \emph{Recall:} say the Malayalam for to tidy away, then the Malayalam for to pull down, then say \emph{kāykkuka} again
\end{itemize}

\begin{tcolorbox}[breakable,colback=teal!4,colframe=teal!35!black,title={Before you move on}]
What does \ml{കായ്ക്കുക} mean? (\textquotedblleft{}To bear fruit\textquotedblright{}.) Say it once more.
\end{tcolorbox}
\section[\texorpdfstring{\ml{പടരുക} (paṭaruka)}{പടരുക (paṭaruka)}]{\ml{പടരുക} (paṭaruka) — to spread (of a creeper, fire, illness)}

\label{lesson:ML-C311-pataruka}



\begin{quote}
Before the new one: say the Malayalam for to pull down, then the Malayalam for to bear fruit.
\end{quote}

\subsection*{You'll want to know: \ml{പടരുക}}
\textbf{\ml{പടരുക}} — \emph{paṭaruka} — \textquotedblleft{}to spread (of a creeper, fire, illness)\textquotedblright{}.

\begin{grammarlens}[title={What you've built}]
One more word for this chapter.
\end{grammarlens}

\subsection*{Your turn}
\begin{itemize}
\raggedright
  \item \emph{Say it:} \emph{paṭaruka}
  \item \emph{Say it:} \emph{paṭaruka}, once more
  \item \emph{Say it:} say \emph{paṭaruka}
  \item \emph{Recall:} say the Malayalam for to pull down, then the Malayalam for to bear fruit, then say \emph{paṭaruka} again
\end{itemize}

\begin{tcolorbox}[breakable,colback=teal!4,colframe=teal!35!black,title={Before you move on}]
What does \ml{പടരുക} mean? (\textquotedblleft{}To spread (of a creeper, fire, illness)\textquotedblright{}.) Say it once more.
\end{tcolorbox}
